\documentclass[tikz]{standalone}
\usepackage{pgfplots}
\usetikzlibrary{decorations.markings, calc, intersections}
\usepgfplotslibrary{fillbetween}

	\usepackage{fontspec}
	\setsansfont{CMU Sans Serif}%{Arial}
	\setmainfont{CMU Serif}%{Times New Roman}
	\setmonofont{CMU Typewriter Text}%{Consolas}
	\defaultfontfeatures{Ligatures={TeX}}
	\usepackage[math-style=TeX]{unicode-math}


\begin{document}

\begin{tikzpicture}[
  midarr/.style={decoration={markings,mark=at position #1 with {\arrow{stealth}}},
  postaction={decorate}}, midarr/.default=0.5, plotred/.style={red!80!black}]

  \def\xmax{10}
  \def\ymax{8}

  % === ОСІ ===
%  \draw[gray!30] (0,0) grid (\xmax, \ymax);
  \draw[->,thick] (0,-0.1*\ymax) -- (0,\ymax) node[anchor=south east] {$P$};
  \draw[->,thick] (-0.1*\xmax,0) -- (\xmax,0) node[anchor=north east] {$V$};

  % === ГОЛОВНИЙ ЦИКЛ ===
  \coordinate (A) at (1.5, 0);
  \coordinate (B) at (1.5, 2);
  \coordinate (C) at (1.5, 6);
  \coordinate (D) at (7, 6);
  \coordinate (E) at (7, 2);
  \coordinate (F) at (7, 0);

  \path (B) -- coordinate[pos=0.25] (1) (C);
  \path (C) -- coordinate[pos=0.75] (2) (D);
  \path (D) -- coordinate[pos=0.25] (3) (E);
  \path (E) -- coordinate[pos=0.75] (4) (B);


%
%  % === МАЛІ ЦИКЛИ КАРНО ===
%  \foreach \x in {1.75, 2.0, ..., 7.25} {
%    \foreach \y in {1.5, 2.5, ..., 5.5} {
%
%      % Задаємо висоту та ширину для кожного ромба
%      \def\height{0.4}  % висота (відповідає висоті ізотерми)
%      \def\width{0.3}   % ширина (відповідає адіабатам)
%
%      % Позиції для ромбів
%      \coordinate (A) at (\x,\y);
%      \coordinate (B) at ($(A) + (\width, -0.5*\height)$);  % Ізотерма
%      \coordinate (C) at ($(B) + (-\width, 1.2)$);          % Адіабата
%      \coordinate (D) at ($(C) + (-\width, \height)$);  % Ізотерма
%%      \coordinate (E) at ($(D) + (-0.6, -1.2)$);        % Адіабата
%
%%    \node at (A) {$A$};
%%\node at (B) {$B$};
%%\node at (C) {$C$};
%%\node at (D) {$D$};
%
%      % Малюємо ромб (замкнуту лінію)
%      \draw[blue!70!black, thick] (A)  to[bend right=5] (B);
%      \draw[blue!70!black, thick] (B) to[bend left=5] (C);
%      \draw[blue!70!black, thick] (C) to[bend left=5] (D);
%      \draw[blue!70!black, thick] (D) to[bend right=5] (A);
%    }
%  }

\begin{scope}[opacity=0.5]
\clip (0,0) rectangle (\xmax, \ymax);
%\clip (1) to[out=90, in=180] (2) to[out=0, in=90] (3) to[out=270, in=0] (4) to[out=180, in=270] (1);
\foreach \t in {3,...,38} {
\ifnum\t>3
\draw[domain=0.5:\xmax, smooth, blue] plot (\x, {\t/\x});
\fi
}
\foreach \t in {3, 3.25, 3.5, ...,15} {
\draw[domain=0.5:\xmax, smooth, red] plot (\x, {\t/\x^(0.5)});
}
\end{scope}

\begin{scope}[]
\clip (0,0) rectangle (\xmax, \ymax);
%\clip (1) to[out=90, in=180] (2) to[out=0, in=90] (3) to[out=270, in=0] (4) to[out=180, in=270] (1);

\foreach \is/\adt/\adb in {
4/3.75/3.00,
5/4.5/3.25,
6/5/3.5,
7/5.5/4,
8/6/4.25,
9/6.5/4.5,
10/7/5,
11/7.5/5.25,
12/8/5.5,
13/8.25/5.75,
14/8.75/6.25,
15/9/6.5,
16/9.5/6.75,
17/9.75/7,
18/10.25/7.5,
19/10.5/7.75,
20/10.75/8,
21/11/8.5,
22/11.5/8.75,
23/11.75/9.25,
24/12.00/9.5,
25/12.25/9.75,
26/12.5/10.25,
27/12.75/10.5,
28/13/11,
29/13.25/11.25,
30/13.5/11.5,
31/13.75/12,
32/14.00/12.25,
33/14.25/12.75,
34/14.5/13,
35/14.5/13.5,
36/14.75/13.75,
37/15.00/14.25
} {

\pgfmathsetmacro\a{\is}
\pgfmathsetmacro\A{\a+1}
\pgfmathsetmacro\b{\adb}
\pgfmathsetmacro\B{\adt}

\pgfmathsetmacro\XaB{(\a/\B)^2}
\pgfmathsetmacro\XAb{(\A/\b)^2}
\pgfmathsetmacro\Xab{(\a/\b)^2}
\pgfmathsetmacro\XAB{(\A/\B)^2}

\draw[domain=\XaB:\Xab, smooth, blue, name path=11] plot (\x, {\a/\x});
\draw[domain=\XAB:\XAb, smooth, blue, name path=12] plot (\x, {\A/\x});

\draw[domain=\XaB:\XAB, smooth, red, name path=8] plot (\x, {\B/\x^(0.5)});
\draw[domain=\Xab:\XAb, smooth, red, name path=5] plot (\x, {\b/\x^(0.5)});

\tikzfillbetween[of=11 and 12]{blue, opacity=0.3};
}

\end{scope}




  \draw[plotred, midarr, ultra thick] (1) to[out=90, in=180] (2);
  \draw[plotred, midarr, ultra thick] (2) to[out=0, in=90] (3);
  \draw[plotred, midarr, ultra thick] (3) to[out=270, in=0] (4);
  \draw[plotred, midarr, ultra thick] (4) to[out=180, in=270] (1);

%  \node[anchor=west] (formula) at (0.5,1) {$\frac{Q_C}{T_C} + \frac{Q_h}{T_h} = 0$}  ;

     \draw[blue] (7, 7.5) -- ++(1,0) node[right] {Ізотерма};
     \draw[red]  (7, 7.0) -- ++(1,0) node[right] {Адіабата};
     \draw[blue!50, line width=5]  (7, 6.5) -- ++(1,0) node[right, text=black] {$\frac{Q_C}{T_C} +
     \frac{Q_h}{T_h} = 0$};

\end{tikzpicture}

\end{document}
